% 数学ノート集・共通テンプレート v2（積分論2026年度第1部の版面を参考）
% A5判。章内で定義・定理・例・注意等の番号を共通化。
% コンパイル： platex -interaction=nonstopmode math_note_template_v2.tex (2回)
%              dvipdfmx math_note_template_v2.dvi
\documentclass[a5paper,10pt,onecolumn]{jsarticle}
\usepackage[dvipdfmx,top=22mm,bottom=20mm,inner=19mm,outer=18mm,headheight=13pt,headsep=6mm]{geometry}
\usepackage{amsmath,amssymb,amsthm,mathtools}
\usepackage{enumitem,fancyhdr,titlesec}
\usepackage[dvipdfmx,bookmarks=true,colorlinks=false]{hyperref}
\usepackage{pxjahyper}

% ----- 差し替える文書情報 -----
\newcommand{\NoteTitle}{Galois理論と定規・コンパス作図}
\newcommand{\NoteSeries}{数学ノート集}

% ----- 文字・組版 -----
% 参考PDFの質素な明朝＋Computer Modernの印象を引き継ぎながら，
% A5判で窮屈にならないよう本文の行間を少し確保する．
\linespread{1.13}
\setlength{\parindent}{1zw}
\setlength{\parskip}{0pt}
\setlength{\abovedisplayskip}{6pt plus 2pt minus 2pt}
\setlength{\belowdisplayskip}{6pt plus 2pt minus 2pt}
\setlength{\abovedisplayshortskip}{3pt plus 2pt minus 1pt}
\setlength{\belowdisplayshortskip}{4pt plus 2pt minus 1pt}
\setlist[enumerate]{topsep=3pt,itemsep=1pt,parsep=0pt,leftmargin=2.4zw}
\titleformat{\section}{\normalfont\normalsize\bfseries}{\thesection}{0.8zw}{}
\titlespacing*{\section}{0pt}{2.0ex plus .4ex}{.9ex}
\titleformat{\subsection}{\normalfont\normalsize\bfseries}{\thesubsection}{0.7zw}{}
\titlespacing*{\subsection}{0pt}{1.5ex plus .3ex}{.6ex}
\raggedbottom

% ----- 柱とページ番号：ページ上部、罫線なし -----
\fancypagestyle{mathnote}{%
  \fancyhf{}%
  \fancyhead[L]{\footnotesize \NoteSeries\ （\NoteTitle）}%
  \fancyhead[R]{\footnotesize \thepage}%
  \renewcommand{\headrulewidth}{0pt}%
  \renewcommand{\footrulewidth}{0pt}%
}
\pagestyle{mathnote}

% ----- 節内通し番号の定義・定理群 -----
\newtheoremstyle{mathnotetheorem}
 {4.5pt}{4.5pt}{\normalfont}{}%
 {\bfseries}{}{0.65zw}%
 {\thmname{#1}\thmnumber{ #2}\thmnote{（#3）}}
\theoremstyle{mathnotetheorem}
\newtheorem{theorem}{定理}[section]
\newtheorem{definition}[theorem]{定義}
\newtheorem{lemma}[theorem]{補題}
\newtheorem{proposition}[theorem]{命題}
\newtheorem{corollary}[theorem]{系}
\newtheorem{example}[theorem]{例}
\newtheorem{remark}[theorem]{注意}
\renewcommand{\proofname}{証明}
\renewcommand{\qedsymbol}{$\square$}
\makeatletter
\renewenvironment{proof}[1][\proofname]{%
  \par\pushQED{\qed}\normalfont\topsep4\p@\@plus2\p@\relax
  \trivlist\item[\hskip\labelsep\bfseries #1]\ignorespaces
}{\popQED\endtrivlist\@endpefalse}
\makeatother

% ----- 使うノートで適宜拡張する記号 -----
\newcommand{\R}{\mathbb{R}}
\newcommand{\norm}[1]{\lVert #1\rVert}
\newcommand{\abs}[1]{\lvert #1\rvert}

\begin{document}
\section*{読了後にできること}
本稿を読むための前提は，定規・コンパス作図，体の二次拡大．
本稿の具体的な到達課題は，作図可能性を体の拡大次数で判定する．

\section{作図可能な点と数}

本稿では平面を複素平面 $\mathbb C$ と同一視し，
初期点を $0,1$ とする．
定規・コンパスによる操作とは，
既知の二点を通る直線，既知の中心と半径で定まる円を描き，
それらの交点を新しい既知の点に加える操作をいう．

\begin{definition}
有限回の定規・コンパス操作で得られる点を
\emph{作図可能点}といい，
その複素数座標を作図可能数という．
作図可能数全体を $\mathcal C\subset\mathbb C$ と書く．
\end{definition}

\begin{lemma}[初等演算の作図可能性]\label{lem:arithmetic}
$a,b\in\mathcal C$ なら，
$a+b$，$a-b$，$ab$ は作図可能である．
$b\ne0$ なら $a/b$ も作図可能である．
さらに $a\in\mathcal C$ の複素共役 $\overline a$ も作図可能である．
\end{lemma}
\begin{proof}
平行線・垂線の作図は定規とコンパスで行える．
この事実を用いて，複素平面の実軸と虚軸への
直交射影を作図できる．
従って作図可能点の実部・虚部の点も作図可能である．
実直線上で平行線を用いた相似三角形の定理により，
長さ $r,s$ の和・差・積・商（$s\ne0$）を作図できる．
実際，原点から二本の半直線を取り，
一方に長さ1と $r$，他方に長さ $s$ を刻み，
対応する二辺を平行にする作図によって
比 $r:s$ と $1:1$ を保つ長さを得る．
符号は実軸の正負方向で表せる．
$a=x+iy$，$b=u+iv$ と書けば
\[
 a+b=(x+u)+i(y+v),\quad
 ab=(xu-yv)+i(xv+yu),
\]
\[
 a/b=\frac{xu+yv}{u^2+v^2}
 +i\frac{yu-xv}{u^2+v^2}
\]
なので，実数の四則演算ができれば
複素数についてもできる．
複素共役は $x+iy\mapsto x-iy$ であるから作図可能である．
\end{proof}

\begin{lemma}[正の平方根と複素平方根]\label{lem:sqrt}
正の作図可能実数 $r$ に対し $\sqrt r$ は作図可能である．
従って任意の $z\in\mathcal C$ に対して
$w^2=z$ となる $w\in\mathcal C$ が存在する．
\end{lemma}
\begin{proof}
実数直線上に長さ $1+r$ の線分を取り，
長さ1と $r$ に分ける点を通る垂線を引く．
この線分を直径とする半円との交点を取れば，
直角三角形の相似性によりその垂線の長さは
幾何平均 $\sqrt r$ に等しい．
従って平方根は作図可能である．

$z=x+iy$ と書くと，$\rho=\sqrt{x^2+y^2}$ は
作図可能である．$y\ne0$ なら
\[
 u=\sqrt{\frac{\rho+x}{2}},\qquad
 v=\operatorname{sgn}(y)\sqrt{\frac{\rho-x}{2}}
\]
を取れば，$u,v$ は実数の作図可能数であり，
$2uv=y$，$u^2-v^2=x$ となるので
$(u+iv)^2=z$ である．
$y=0$ なら $x\ge0$ のとき $w=\sqrt x$，
$x<0$ のとき $w=i\sqrt{-x}$ とすればよい．
\end{proof}

\section{作図操作の代数化}

\begin{lemma}[直線と円の交点]\label{lem:line-circle}
$F\subset\mathbb R$ を四則演算に閉じた体とする．
係数が $F$ に属する二直線，直線と円，
または二円の交点の座標は，
ある二次以下の体拡大 $E/F$ に属する．
\end{lemma}
\begin{proof}
二直線の交点は二元一次連立方程式を解くことで得られるため，
座標は $F$ に属する．
直線と円の場合，
直線を $(x,y)=(x_0,y_0)+t(u,v)$ と表示でき，
これを円の方程式
$(x-a)^2+(y-b)^2=r^2$ に代入すると
$t$ に関する二次方程式が得られる．
その係数は $F$ に属するので，
交点の座標は判別式の平方根を加えた体に属する．
二円の場合には両円の方程式を引けば，
中心が異なるなら一次方程式が得られ，
直線と円の場合に帰着する．
中心が同じなら二円は一致するか交わらないので，
新しい孤立した交点は得られない．
従ってどの場合も高々二次拡大で足りる．
\end{proof}

\begin{theorem}[作図可能性の体論的特徴づけ]\label{thm:constructible}
複素数 $z$ が作図可能であるための必要十分条件は，
有限個の体からなる塔
\[
 \mathbb Q=F_0\subset F_1\subset\cdots\subset F_r
 \subset\mathbb C,\qquad [F_i:F_{i-1}]\le2
\]
が存在し，$z\in F_r$ となることである．
\end{theorem}
\begin{proof}
作図可能点の座標は最初 $0,1$ から始まり，
線・円の方程式の係数は既知点の実部・虚部から
四則演算で得られる．
補題\ref{lem:line-circle}により，
一回の作図で新しい座標を付け加える操作は
高々二次の拡大で実現できる．
複数の新しい座標が現れても有限個の二次拡大に分割できる．
従って作図可能数は所望の塔に含まれる．

逆に $F_{i}=F_{i-1}(\alpha)$ が二次拡大なら，
$\alpha$ は
$\alpha^2+a\alpha+b=0$（$a,b\in F_{i-1}$）
を満たすから，二次方程式の公式によって
\[
 \alpha=\frac{-a\pm\sqrt{a^2-4b}}2
\]
と書ける．
塔の最初の体 $\mathbb Q$ の元は四則演算で作図可能であり，
補題\ref{lem:arithmetic}と補題\ref{lem:sqrt}により，
各段階で体を構成する全ての元が作図可能となる．
従って最後の体 $F_r$ の元 $z$ も作図可能である．
\end{proof}

\begin{corollary}[次数による障害]\label{cor:degree}
$z\in\mathbb C$ が作図可能で，$z$ が $\mathbb Q$ 上代数的なら，
\[
 [\mathbb Q(z):\mathbb Q]=2^k
\]
となる非負整数 $k$ が存在する．
\end{corollary}
\begin{proof}
定理\ref{thm:constructible}の塔について
$[F_r:\mathbb Q]$ は2の冪である．
$\mathbb Q\subset\mathbb Q(z)\subset F_r$ なので
塔の定理から $[\mathbb Q(z):\mathbb Q]$ は
$[F_r:\mathbb Q]$ の約数である．
2の冪の正の約数は2の冪なので結論を得る．
\end{proof}

\section{正規拡大と2群}

\begin{lemma}[有限2群の正規列]\label{lem:2group}
有限群 $G$ の位数が2の冪なら，
\[
 G=G_0\triangleright G_1\triangleright\cdots
 \triangleright G_r=1
\]
で，各 $[G_i:G_{i+1}]=2$ となる列が存在する．
\end{lemma}
\begin{proof}
$G\ne1$ なら有限 $p$ 群の中心は非自明である．
これは類等式
$|G|=|Z(G)|+\sum_i[G:C_G(g_i)]$ において
中心外の各共役類の大きさが2の倍数であることから従う．
従って $Z(G)$ は非自明な2群であり，
Cauchyの定理によって位数2の元 $z$ を持つ．
$\langle z\rangle$ は中心に含まれるため正規部分群であり，
商 $G/\langle z\rangle$ も2群である．
位数について帰納法を適用すると，
商群が指数2の正規列を持つ．
その逆像を取って $G$ の正規列へ引き戻せば，
各指数は2であり，最後に $\langle z\rangle\triangleright1$
を加えて所望の列を得る．
\end{proof}

\begin{theorem}[Galois拡大における作図可能性]
\label{thm:galois-2}
$L/\mathbb Q$ を有限Galois拡大とする．
$L$ の全ての元が作図可能であるための必要十分条件は
$\operatorname{Gal}(L/\mathbb Q)$ が2群であることである．
\end{theorem}
\begin{proof}
$L$ の全ての元が作図可能なら，
特に原始元定理により
$L=\mathbb Q(\alpha)$ となる原始元 $\alpha$ は作図可能である．
系\ref{cor:degree}から
$[L:\mathbb Q]=[\mathbb Q(\alpha):\mathbb Q]$
は2の冪である．
$L/\mathbb Q$ はGaloisなので
$|\operatorname{Gal}(L/\mathbb Q)|=[L:\mathbb Q]$
であり，Galois群は2群となる．

逆にGalois群が2群なら補題\ref{lem:2group}によって
指数2の正規列を持つ．
Galois対応を適用すると
\[
 \mathbb Q=F_0\subset F_1\subset\cdots\subset F_r=L
\]
で $[F_i:F_{i-1}]=2$ となる体の塔が得られる．
定理\ref{thm:constructible}によって
$L$ の全ての元は作図可能である．
\end{proof}

\section{古典的作図不可能性}

\begin{theorem}[立方体倍積問題]\label{thm:duplication}
単位立方体と体積が二倍の立方体の一辺の長さ
$\sqrt[3]2$ は定規・コンパスで作図できない．
\end{theorem}
\begin{proof}
$X^3-2$ は素数2に関するEisensteinの判定法により
$\mathbb Q$ 上既約である．
従って
$[\mathbb Q(\sqrt[3]2):\mathbb Q]=3$
である．
これは2の冪ではないから
系\ref{cor:degree}によって作図不可能である．
\end{proof}

\begin{theorem}[一般の角の三等分は不可能]
\label{thm:trisection}
$60^\circ$ の角を定規・コンパスで三等分することはできない．
\end{theorem}
\begin{proof}
$60^\circ$ を三等分できるなら $20^\circ$ を作図でき，
$x=2\cos20^\circ$ も作図可能となる．
三倍角の公式から
\[
 2\cos60^\circ
 =(2\cos20^\circ)^3-3(2\cos20^\circ)
\]
なので
\[
 x^3-3x-1=0.
\]
有理根定理からこのモニック三次式の有理根候補は
$\pm1$ だけであり，
$f(1)=-3$，$f(-1)=1$ なので有理根はない．
三次式は有理根を持たなければ既約だから
$[\mathbb Q(x):\mathbb Q]=3$ である．
系\ref{cor:degree}に反するため作図不可能である．
\end{proof}

\section{円分多項式の既約性}

\begin{definition}
$n\ge1$ に対し，原始 $n$ 乗根の集合を
$\mu_n^{\mathrm{prim}}$ と書き，
\[
 \Phi_n(X)=\prod_{\zeta\in\mu_n^{\mathrm{prim}}}(X-\zeta)
\]
を\emph{第 $n$ 円分多項式}という．
\end{definition}

\begin{lemma}[円分多項式は整数係数]\label{lem:cyclotomic-integral}
任意の $n\ge1$ について
$\Phi_n(X)\in\mathbb Z[X]$ が成り立ち，
\[
 X^n-1=\prod_{d\mid n}\Phi_d(X).
\]
\end{lemma}
\begin{proof}
$X^n-1$ の根は位数 $d\mid n$ によって一意に分類されるので，
複素数係数多項式として積公式が従う．
$n=1$ では $\Phi_1(X)=X-1$ であり整数係数である．
$n>1$ について帰納法で
$d<n$ の $\Phi_d$ が整数係数であるとする．
\[
 \Phi_n(X)=\frac{X^n-1}{\prod_{\substack{d\mid n\\d<n}}\Phi_d(X)}
\]
であり，分母はモニックな整数係数多項式である．
モニック多項式での割り算を整数係数環で実行すると，
商と余りはともに整数係数となる．
複素数係数の積公式から余りは零なので，
商 $\Phi_n$ は整数係数である．
\end{proof}

\begin{theorem}[円分多項式の既約性]\label{thm:cyclotomic-irred}
$\Phi_n(X)$ は $\mathbb Q$ 上既約である．
従って
\[
 [\mathbb Q(\zeta_n):\mathbb Q]=\varphi(n)
\]
が成立する．
\end{theorem}
\begin{proof}
$\zeta$ を原始 $n$ 乗根とし，
$f\in\mathbb Z[X]$ を $\zeta$ のモニック最小多項式とする．
Gaussの補題と補題\ref{lem:cyclotomic-integral}から
$f$ は $\Phi_n$ のモニックな整数係数因子である．
任意の素数 $p\nmid n$ に対して
$\zeta^p$ も $f$ の根であることを示す．
もしそうでなければ，$\zeta^p$ の最小多項式
$g\in\mathbb Z[X]$ は $f$ と異なり，
$f$ と $g$ はともに $X^n-1$ の異なる既約因子となる．
$g(\zeta^p)=0$ だから
$f$ は $g(X^p)$ を割る．
係数を $\mathbb F_p$ に移すと
\[
 g(X^p)\equiv g(X)^p\pmod p
\]
である．
従って $\overline f$ と $\overline g$ は
$\mathbb F_p[X]$ で共通の非定数因子を持つ．
一方 $fg$ は $X^n-1$ を $\mathbb Z[X]$ で割るので，
$\overline f\,\overline g$ は
$X^n-1$ を $\mathbb F_p[X]$ で割る．
$p\nmid n$ だから
$(X^n-1)'=nX^{n-1}$ は $X^n-1$ と互いに素である．
従って $X^n-1$ は $\mathbb F_p$ 上平方因子を持たず，
$\overline f,\overline g$ が共通因子を持つことに矛盾する．
よって $\zeta^p$ も $f$ の根である．

任意の $a$ が $\gcd(a,n)=1$ を満たすなら，
$a$ の素因数は全て $n$ と互いに素である．
上の結論を $\zeta$ の各共役に繰り返し適用することで
$\zeta^a$ も $f$ の根となる．
従って $f$ は全ての原始 $n$ 乗根を根として含む．
$\Phi_n$ はちょうどそれらを根とするモニック多項式であり，
$f\mid\Phi_n$ だから $f=\Phi_n$ である．
$\Phi_n$ の次数は原始 $n$ 乗根の個数
$\varphi(n)$ であるため，後半も従う．
\end{proof}

\section{Gauss--Wantzelの定理}

\begin{theorem}[正多角形の作図可能性]\label{thm:gauss-wantzel}
$n\ge3$ とする．正 $n$ 角形が定規・コンパスで作図可能である
ための必要十分条件は
\[
 n=2^a p_1\cdots p_r
\]
と書けることである．
ここで $a\ge0$，$p_1,\dots,p_r$ は相異なる
Fermat素数，すなわち
$p_j=2^{2^{k_j}}+1$ となる素数である．
\end{theorem}
\begin{proof}
正 $n$ 角形が作図可能であることは，
原始 $n$ 乗根 $\zeta_n$ が作図可能であることと同値である．
$\mathbb Q(\zeta_n)/\mathbb Q$ は円分多項式の分解体で
有限Galois拡大である．
定理\ref{thm:cyclotomic-irred}から次数は
$\varphi(n)$ である．
定理\ref{thm:galois-2}により，
正 $n$ 角形が作図可能であることと
$\varphi(n)$ が2の冪であることは同値である．

$n=2^a\prod_{j=1}^r p_j^{e_j}$ と
相異なる奇素数 $p_j$ によって素因数分解する．
Euler関数の乗法性と素数冪上の公式から
\[
 \varphi(n)=\varphi(2^a)
 \prod_{j=1}^r p_j^{e_j-1}(p_j-1).
\]
これが2の冪となるためには，各 $e_j=1$ かつ
$p_j-1$ が2の冪でなければならない．
逆にこれらの条件があれば明らかに
$\varphi(n)$ は2の冪となる．
$p_j=2^m+1$ が素数なら $m$ は2の冪であることを示す．
もし $m=2^kr$ と書け，$r>1$ が奇数なら，
$2^m+1=(2^{2^k})^r+1$ は
$2^{2^k}+1$ を因数に持つため合成数となる．
従って $m=2^k$ である．
ゆえに条件は表示したFermat素数による形と同値である．
\end{proof}

\begin{example}
正17角形は作図可能である．
実際 $17=2^{2^2}+1$ はFermat素数であり，
$\varphi(17)=16$ は2の冪である．
一方，正7角形は作図できない．
\end{example}
\begin{proof}
$17$ が素数であることは
$2,3,4$ のうち素数2,3で割り切れないことから従う．
定理\ref{thm:gauss-wantzel}を適用して正17角形の
作図可能性を得る．
正7角形については
$\varphi(7)=6$ が2の冪ではないため
同じ定理から作図不可能である．
\end{proof}

\begin{remark}
作図可能性の定理は，
幾何学的操作を二次体拡大の塔へ翻訳し，
有限Galois群の2群性で特徴づけることで成立する．
「次数が2の冪」という必要条件は一般の代数数について
十分条件ではないが，円分拡大のようにGaloisである場合には
Galois対応によって十分条件となる点が重要である．
\end{remark}

\section{到達課題の完全解答}
\noindent\textbf{正三角形と立方体倍積の比較．}
\begin{proposition}
$\sqrt3$ は $\mathbb Q$ 上作図可能であるが，
$\sqrt[3]2$ は作図不可能である．
\end{proposition}
\begin{proof}
$\sqrt3$ は $x^2-3$ の根であり，
平方根を取る操作は定規・コンパスで実行できる．
一方 $x^3-2$ はEisensteinの判定法で既約で，
$[\mathbb Q(\sqrt[3]2):\mathbb Q]=3$ である．
作図可能な実数は $\mathbb Q$ からの
二次拡大の塔に含まれるため，
その最小多項式の次数は2の冪を割り切る必要がある．
3は2の冪ではないので $\sqrt[3]2$ は作図不可能である．
\end{proof}
\noindent\textbf{到達確認．}
$\sqrt{2+\sqrt2}$ は作図可能か．
\par\noindent\textbf{解答．}
$\mathbb Q\subset\mathbb Q(\sqrt2)
\subset\mathbb Q(\sqrt2,\sqrt{2+\sqrt2})$
は各段階で次数高々2の拡大なので作図可能である．
\end{document}